\chapter{Direito Constitucional}

\section{Constituição, constitucionalismo e supremacia constitucional}

A Constituição é a norma fundamental que organiza o Estado, distribui competências, estrutura os Poderes, reconhece direitos e impõe limites ao exercício do poder político. No Brasil, a Constituição de 1988 ocupa posição superior no sistema normativo e serve como parâmetro de validade para leis e atos estatais.

\subsection{Constitucionalismo}

Constitucionalismo é o movimento histórico e jurídico de limitação do poder por normas superiores, com proteção de direitos e organização institucional. A ideia moderna de Constituição associa-se, portanto, a dois grandes eixos: \textbf{organização do poder} e \textbf{garantia de direitos}.

\subsection{Supremacia formal e material}

A supremacia \termo{formal} decorre da rigidez constitucional: alterar a Constituição exige procedimento mais solene do que editar legislação ordinária. A supremacia \termo{material} relaciona-se à centralidade de seus conteúdos para a organização do Estado e proteção de direitos.

O controle de constitucionalidade é consequência da supremacia: norma inferior incompatível com a Constituição não pode permanecer válida no sistema nos termos previstos pelo ordenamento.

\section{Classificação da Constituição de 1988}

Na classificação doutrinária tradicional, a CF/1988 é:
\begin{itemize}
\item \textbf{escrita}: consolidada em documento solene;
\item \textbf{promulgada}: elaborada por assembleia constituinte representativa;
\item \textbf{dogmática}: construída de forma sistematizada a partir de ideias políticas e jurídicas;
\item \textbf{rígida}: sua alteração exige procedimento mais difícil que o legislativo ordinário;
\item \textbf{analítica}: disciplina extensa variedade de matérias;
\item \textbf{formal}: considera constitucionais as normas inseridas no texto segundo o procedimento constituinte, independentemente de seu conteúdo material.
\end{itemize}

\section{Poder constituinte}

\subsection{Poder constituinte originário}

É o poder responsável por criar nova ordem constitucional. A doutrina o caracteriza como inicial, autônomo e juridicamente ilimitado em relação à ordem anterior, embora sujeito a debates sobre limites de natureza política, histórica ou suprapositiva.

\subsection{Poder constituinte derivado reformador}

É exercido por emendas constitucionais e encontra limites na própria Constituição.

Há limites:
\begin{itemize}
\item \textbf{formais}: iniciativa, quórum e procedimento;
\item \textbf{circunstanciais}: situações em que não se pode emendar a Constituição;
\item \textbf{materiais}: cláusulas pétreas;
\item \textbf{implícitos}: construídos a partir da lógica do próprio sistema constitucional.
\end{itemize}

A proposta de emenda é aprovada em cada Casa do Congresso Nacional, em dois turnos, por três quintos dos votos dos respectivos membros.

\subsection{Cláusulas pétreas}

Não será objeto de deliberação proposta tendente a abolir:
\begin{itemize}
\item forma federativa de Estado;
\item voto direto, secreto, universal e periódico;
\item separação dos Poderes;
\item direitos e garantias individuais.
\end{itemize}

A proteção recai também contra propostas \emph{tendentes} a abolir, não apenas contra supressão textual expressa.

\section{Aplicabilidade das normas constitucionais}

Uma classificação doutrinária muito cobrada distingue normas de eficácia plena, contida e limitada.

\subsection{Eficácia plena}

Produzem seus efeitos essenciais desde a entrada em vigor da Constituição, sem depender de lei integrativa para serem aplicáveis em seu núcleo.

\subsection{Eficácia contida}

Também possuem aplicabilidade imediata, mas seu alcance pode ser restringido nos termos autorizados pela própria Constituição.

\subsection{Eficácia limitada}

Dependem de atuação normativa ou institucional posterior para produzir integralmente os efeitos pretendidos. Ainda assim, não são juridicamente inúteis: vinculam o legislador, condicionam interpretação e impedem atuação estatal incompatível com o programa constitucional.

\section{Princípios fundamentais}

Os arts. 1º a 4º da Constituição organizam fundamentos da República, objetivos fundamentais e princípios das relações internacionais.

\subsection{Fundamentos}

São fundamentos da República Federativa do Brasil: soberania, cidadania, dignidade da pessoa humana, valores sociais do trabalho e da livre iniciativa e pluralismo político.

\subsection{Objetivos fundamentais}

Entre os objetivos estão construir sociedade livre, justa e solidária; garantir desenvolvimento nacional; erradicar pobreza e marginalização e reduzir desigualdades; promover o bem de todos sem preconceitos e discriminações.

\begin{atencao}
Fundamentos, objetivos e princípios das relações internacionais formam listas distintas. A banca frequentemente desloca um item de uma categoria para outra.
\end{atencao}

\section{Direitos e garantias fundamentais}

Direitos fundamentais são posições jurídicas protegidas constitucionalmente. Garantias são instrumentos e mecanismos destinados a assegurar esses direitos.

A aplicação imediata prevista no art. 5º, §1º, não significa que todo direito produza qualquer efeito independentemente de regulamentação. Também não significa que direitos fundamentais sejam absolutos.

\subsection{Igualdade}

Igualdade possui dimensão formal — tratamento igual perante a lei — e dimensão material, que admite diferenciações juridicamente justificadas para reduzir desigualdades ou tratar situações distintas de modo proporcional.

\subsection{Legalidade}

No âmbito dos particulares, ninguém é obrigado a fazer ou deixar de fazer algo senão em virtude de lei. Na Administração, a legalidade possui feição mais estrita, vinculada às competências e finalidades definidas pelo ordenamento.

\subsection{Intimidade, vida privada, honra e imagem}

São bens constitucionalmente protegidos e podem entrar em tensão com liberdade de expressão, transparência e interesse público. A solução exige ponderação e aplicação das regras específicas do caso.

\subsection{Devido processo, contraditório e ampla defesa}

O devido processo legal protege contra privação arbitrária de liberdade ou bens. Contraditório e ampla defesa alcançam processos judiciais e administrativos nos termos constitucionais.

Contraditório não é apenas conhecimento formal; pressupõe possibilidade de participação e influência na formação da decisão quando juridicamente cabível.

\section{Remédios constitucionais}

\begin{table}[ht]
\centering
\small
\begin{tabularx}{\textwidth}{l X X}
\toprule
Instrumento & Finalidade central & Observação \\
\midrule
Habeas corpus & proteger liberdade de locomoção & cabe contra ilegalidade ou abuso que ameace ou restrinja ir e vir. \\
Mandado de segurança & proteger direito líquido e certo & não amparado por HC ou HD, diante de ilegalidade ou abuso de poder. \\
Habeas data & acesso/retificação de dados pessoais & nos registros e hipóteses previstas constitucional e legalmente. \\
Mandado de injunção & enfrentar omissão regulamentadora & quando falta de norma inviabiliza exercício de direito, liberdade ou prerrogativa constitucional. \\
Ação popular & combater ato lesivo a bens constitucionais protegidos & legitimidade ativa do cidadão. \\
\bottomrule
\end{tabularx}
\end{table}

\section{Direitos sociais}

Direitos sociais buscam assegurar condições materiais de dignidade e participação social. A Constituição inclui educação, saúde, alimentação, trabalho, moradia, transporte, lazer, segurança, previdência social, proteção à maternidade e à infância e assistência aos desamparados, entre outros desdobramentos constitucionais.

A efetivação de direitos sociais envolve legislação, políticas públicas, orçamento e controle judicial em situações constitucionalmente relevantes. Não se pode reduzir o tema à afirmação de que todo direito social depende exclusivamente de escolha política discricionária.

\section{Nacionalidade}

A Constituição distingue brasileiros natos e naturalizados. Algumas funções são reservadas a brasileiros natos. Nacionalidade é vínculo jurídico-político com o Estado e não se confunde com cidadania, que se relaciona especialmente ao exercício de direitos políticos.

\section{Direitos políticos}

Direitos políticos disciplinam participação do cidadão na formação da vontade estatal. Incluem sufrágio, voto, alistamento, elegibilidade e regras de perda ou suspensão.

\subsection{Voto}

O voto é direto e secreto, com valor igual para todos. A obrigatoriedade ou facultatividade depende das categorias constitucionais de idade e condição.

\subsection{Perda e suspensão}

A Constituição veda cassação de direitos políticos e prevê hipóteses específicas de perda ou suspensão. Em prova, a classificação correta depende da redação constitucional vigente e da situação apresentada.

\section{Organização do Estado e federação}

A República Federativa do Brasil é composta por União, Estados, Distrito Federal e Municípios, todos autônomos nos termos constitucionais. Soberania pertence à República Federativa do Brasil como Estado federal, não a cada ente isoladamente.

\subsection{Autonomia}

Autonomia envolve capacidade de auto-organização, autogoverno, autoadministração e autolegislação dentro das competências constitucionais.

\section{Repartição de competências}

A Constituição combina diferentes técnicas de repartição.

\subsection{Competências materiais}

Competências materiais dizem respeito à execução de atividades administrativas. Algumas são exclusivas da União; outras são comuns aos entes federativos.

\subsection{Competências legislativas}

A União possui competências legislativas privativas. Há também competência concorrente entre União, Estados e Distrito Federal em matérias definidas constitucionalmente. Municípios legislam sobre interesse local e suplementam legislação federal e estadual no que couber.

Na competência concorrente, a União estabelece normas gerais. Estados podem exercer competência suplementar e, na ausência de norma geral federal, competência legislativa plena para atender peculiaridades, sujeita aos efeitos constitucionais de posterior lei federal.

\begin{exemplo}[norma geral superveniente]
Se um Estado edita disciplina plena em matéria concorrente por ausência de norma geral federal e, posteriormente, a União edita norma geral, a lei estadual não é simplesmente revogada em bloco. Sua eficácia fica suspensa no que for contrário à norma geral federal, conforme o regime constitucional.
\end{exemplo}

\section{Intervenção}

Intervenção é medida excepcional que limita temporariamente autonomia federativa para enfrentar hipóteses constitucionalmente previstas. Intervenção federal e estadual seguem fundamentos e procedimentos próprios.

A regra é autonomia; intervenção é exceção e deve ser interpretada de forma estrita.

\section{Administração Pública na Constituição}

O art. 37 e seguintes disciplinam princípios, concurso, acumulação, teto, servidores e responsabilidade civil.

\subsection{Concurso público}

Investidura em cargo ou emprego depende, em regra, de concurso de provas ou de provas e títulos, ressalvadas nomeações para cargo em comissão declarado em lei de livre nomeação e exoneração.

O prazo de validade do concurso é de até dois anos, prorrogável uma vez por igual período.

\subsection{Acumulação remunerada}

A regra é vedação, com exceções constitucionais, condicionadas à compatibilidade de horários e observância dos demais requisitos.

\subsection{Estabilidade}

Servidor nomeado para cargo efetivo em virtude de concurso adquire estabilidade após três anos de efetivo exercício e avaliação especial de desempenho, observadas as regras constitucionais. Estabilidade não torna o cargo absolutamente indisponível: a Constituição prevê hipóteses de perda do cargo.

\section{Separação dos Poderes}

Legislativo, Executivo e Judiciário são independentes e harmônicos. A separação é acompanhada por mecanismos de freios e contrapesos.

Cada Poder exerce função \textbf{típica} e também pode desempenhar funções \textbf{atípicas}. O Legislativo, por exemplo, administra sua estrutura; o Executivo edita atos normativos; o Judiciário administra seus servidores e orçamento.

\section{Poder Legislativo}

No plano federal, o Congresso Nacional é bicameral: Câmara dos Deputados e Senado Federal.

\subsection{Processo legislativo}

O art. 59 prevê elaboração de:
\begin{itemize}
\item emendas à Constituição;
\item leis complementares;
\item leis ordinárias;
\item leis delegadas;
\item medidas provisórias;
\item decretos legislativos;
\item resoluções.
\end{itemize}

\subsection{Lei complementar e lei ordinária}

A lei complementar exige maioria absoluta e é reservada a matérias que a Constituição expressamente lhe atribui. Lei ordinária utiliza, em regra, quórum de maioria simples nas condições constitucionais.

Não existe hierarquia geral e abstrata segundo a qual toda lei complementar seja superior a toda lei ordinária. A distinção decorre principalmente de competência material e procedimento.

\subsection{Medida provisória}

Medidas provisórias são editadas pelo Presidente da República em caso de relevância e urgência e produzem força de lei, sujeitas a limitações materiais, prazo e apreciação pelo Congresso.

\section{Fiscalização contábil, financeira e orçamentária}

O art. 70 determina fiscalização contábil, financeira, orçamentária, operacional e patrimonial quanto a legalidade, legitimidade, economicidade, aplicação de subvenções e renúncia de receitas.

O controle externo federal é exercido pelo Congresso Nacional com auxílio do Tribunal de Contas da União. Cada Poder também mantém sistema de controle interno.

\subsection{Dever de prestar contas}

Deve prestar contas qualquer pessoa física ou jurídica, pública ou privada, que utilize, arrecade, guarde, gerencie ou administre dinheiros, bens ou valores públicos ou pelos quais a União responda, bem como quem assuma obrigação pecuniária em nome dela.

\section{Tribunais de Contas}

Tribunais de Contas não integram o Poder Judiciário. Também não são departamentos subordinados hierarquicamente ao Legislativo. Exercem competências constitucionais próprias no auxílio ao controle externo.

\subsection{Apreciar e julgar contas}

O TCU \textbf{aprecia} as contas prestadas anualmente pelo Presidente da República e emite parecer prévio. Em contraste, \textbf{julga} as contas dos administradores e demais responsáveis por dinheiros, bens e valores públicos e daqueles que derem causa a perda, extravio ou irregularidade de que resulte prejuízo ao erário.

Essa distinção entre \emph{apreciar} e \emph{julgar} é central em concursos de controle.

\subsection{Outras competências}

Entre outras funções, o TCU pode:
\begin{itemize}
\item apreciar, para fins de registro, legalidade de determinados atos de admissão e concessão;
\item realizar inspeções e auditorias;
\item fiscalizar recursos repassados;
\item aplicar sanções previstas em lei;
\item assinar prazo para cumprimento da lei;
\item sustar atos impugnados e comunicar ao Congresso, nos termos constitucionais;
\item representar ao poder competente sobre irregularidades ou abusos.
\end{itemize}

\begin{cebraspe}
``O Tribunal de Contas julga as contas do chefe do Poder Executivo'' é formulação incorreta em regra constitucional federal: o TCU aprecia as contas do Presidente e emite parecer prévio; o julgamento político cabe ao Congresso Nacional. Já as contas de administradores e responsáveis são julgadas pelo Tribunal de Contas.
\end{cebraspe}

\section{Poder Executivo}

O Presidente da República exerce chefia de Estado e chefia de Governo. A Constituição define processo eleitoral, substituição, sucessão, atribuições e responsabilização.

Chefias de Estado envolvem representação da República e relações institucionais externas. Chefias de Governo relacionam-se à direção política e administrativa interna.

\section{Poder Judiciário}

A Constituição organiza STF, CNJ, STJ, Justiça Federal, do Trabalho, Eleitoral, Militar e Justiças estaduais, entre outros órgãos constitucionalmente previstos.

\subsection{Garantias da magistratura}

Vitaliciedade, inamovibilidade e irredutibilidade de subsídio protegem independência funcional, nos termos constitucionais.

\subsection{CNJ}

O Conselho Nacional de Justiça exerce controle da atuação administrativa e financeira do Judiciário e do cumprimento de deveres funcionais dos magistrados. Não funciona como tribunal de recursos destinado a revisar o mérito jurisdicional de decisões judiciais.

\section{Funções essenciais à Justiça}

Ministério Público, Advocacia Pública, Advocacia e Defensoria Pública possuem regime constitucional próprio.

O Ministério Público é instituição permanente, essencial à função jurisdicional, incumbida da defesa da ordem jurídica, do regime democrático e dos interesses sociais e individuais indisponíveis.

\section{Controle de constitucionalidade}

Controle de constitucionalidade verifica compatibilidade entre normas e Constituição.

\subsection{Controle preventivo e repressivo}

Preventivo ocorre antes de a norma ingressar definitivamente no ordenamento, como análise legislativa de constitucionalidade e veto jurídico. Repressivo atua sobre norma já formada, sobretudo pelo Judiciário.

\subsection{Controle difuso}

No controle difuso, a questão constitucional aparece em caso concreto e pode ser examinada pelos órgãos jurisdicionais competentes. A inconstitucionalidade surge como questão prejudicial à solução da controvérsia.

\subsection{Controle concentrado}

No plano federal, o STF julga ações como ADI, ADC, ADO e ADPF, cada qual com objeto, legitimidade e finalidade próprios.

\begin{table}[ht]
\centering
\small
\begin{tabularx}{\textwidth}{l X}
\toprule
Ação & Função central \\
\midrule
ADI & questionar constitucionalidade de lei ou ato normativo federal ou estadual, nos termos constitucionais. \\
ADC & obter declaração de constitucionalidade de lei ou ato normativo federal. \\
ADO & enfrentar omissão inconstitucional. \\
ADPF & evitar ou reparar lesão a preceito fundamental, observada subsidiariedade e regime legal. \\
\bottomrule
\end{tabularx}
\end{table}

\subsection{Efeitos e modulação}

Não é suficiente decorar ``difuso = inter partes'' e ``concentrado = erga omnes''. O sistema brasileiro combina regras de efeitos, precedentes, repercussão geral, súmulas vinculantes e modulação temporal. A resposta depende do instrumento e do regime aplicável.

\section{Finanças públicas e orçamento}

A Constituição estrutura planejamento e orçamento por PPA, LDO e LOA.

\begin{itemize}
\item \textbf{PPA}: estabelece diretrizes, objetivos e metas para período plurianual;
\item \textbf{LDO}: orienta elaboração da LOA e conecta planejamento e orçamento;
\item \textbf{LOA}: estima receitas e fixa despesas para o exercício.
\end{itemize}

Controle orçamentário deve ser compreendido em conjunto com fiscalização de legalidade, legitimidade, economicidade e resultados.

\section{Síntese conceitual}

\mapafuncional{Direito Constitucional}
{Constituição = fundamento de validade + organização do poder + proteção de direitos}
{Teoria constitucional $\rightarrow$ supremacia, poder constituinte e eficácia.\\Direitos fundamentais $\rightarrow$ garantias e remédios.\\Federação $\rightarrow$ autonomia e competências.\\Poderes $\rightarrow$ Legislativo, Executivo, Judiciário e funções essenciais.\\Controle $\rightarrow$ fiscalização financeira + constitucionalidade.}
{Norma constitucional: identifique eficácia.\\Competência: identifique ente e natureza.\\Tribunal de Contas: diferencie apreciar de julgar.\\Processo legislativo: verifique matéria, quórum e iniciativa.\\Controle de constitucionalidade: identifique objeto e instrumento.}
{Autonomia $\neq$ soberania.\\Direito fundamental $\neq$ absoluto.\\TC $\neq$ Judiciário.\\Lei complementar $\neq$ superior a toda lei ordinária.\\Aplicabilidade imediata $\neq$ ausência de regulamentação em qualquer hipótese.}
{Qual é o dispositivo constitucional relevante?\\Quem tem competência?\\A norma é plena, contida ou limitada?\\O ato é função típica ou atípica?\\O Tribunal aprecia ou julga?\\Qual ação de controle é adequada?}

\begin{resumo}
O estudo constitucional precisa conectar teoria e organização institucional. Supremacia explica o controle de constitucionalidade; federação explica a repartição de competências; separação dos Poderes explica controles recíprocos; e os arts. 70 e 71 dão a base constitucional do controle externo que estrutura a atuação dos Tribunais de Contas.
\end{resumo}
